\section{Discussion}\label{sec:discussion}

\subsection{What the agent's internals can and cannot replace}
The results support a precise version of the claim that motivated this study. The forward pass that produces an
LLM agent's proposal carries information about whether that proposal will pass simulation. The information is
available before the simulator runs, costs a fixed amount, and requires no model of the plant. How much
verification it can replace depends on the kind of error the agent makes. When errors are failures to consult
information the agent was given, as in path planning, the internals support certified unchecked execution. They
add to observables even under a flexible probe, and they beat any observable baseline that does not itself compute
part of the validator's answer. When errors are reasoning errors over correctly consulted information, as in CSTR
fault recovery, the internals still identify many failing proposals and reject them with 95--98\% precision. They
cannot, however, certify a proposal as safe, and a flexible surrogate of the plant's response captures what they
add. Region-based grounding makes
this distinction measurable. Its AUROC on development data is a cheap diagnostic of whether internal signals are
worth the extra feature pass for a given task and model.

The asymmetry between rejection and acceptance has a practical reading. A rejection by an internals-only screen
is cheap to trust: it can at worst cost a regeneration, never an unsafe action. It therefore requires no plant
model and no certification for safety. An acceptance is where the guarantee is needed. Our results suggest that,
in physics-heavy tasks, acceptance should rely on features that describe the plant's response, while the agent's
internals contribute mainly by filtering out proposals that are clearly wrong.

\subsection{Plant surrogates versus model internals}
A natural objection is that a screen on plant measurements and proposed setpoints is a learned surrogate of the
digital twin rather than a statement about the agent. The sensitivity analysis confirms this reading: given a
flexible probe, the CSTR observable screen improved by up to 40 points and absorbed every gain from internals
(Fig.~\ref{fig:heatmap}). Where a cheap surrogate of the plant can be learned, it is the better tool for
\emph{accepting} proposals. The study shows what the internals offer beyond it in three ways. First, the
plant-readings-only baseline, which sees the plant state but not the proposed action, skipped nearly as many
calls as the full observable screen for three models, but almost none for Qwen2.5-3B (1.8\%) and Llama-3.2-3B
(4.0\%). For those two models the outcome depends on the agent's action, not only on the fault state. Second, the
internals-only screen uses no plant information at all. With the first-proposal-only rule it avoided 29--65\% of
validator calls while executing at most two failing actions in 2000 closed-loop episodes. Third, a surrogate
must be retrained when the plant changes, whereas the internal signals are tied to the agent. Our shift analysis
shows, however, that this does not make them more robust to a change of plant population. Both kinds of screen
must be recalibrated (Section~\ref{sec:res-shift}).

\subsection{A deployment guide}
Figure~\ref{fig:guide} condenses the results into a procedure for engineers who deploy an LLM agent behind a
simulator. Three design choices mattered more than the choice of signal.
\begin{enumerate}
\item \emph{Thresholds need a confidence bound.} A point-estimate accept threshold overshot the tolerance on
certification data. Even with the bound-based rule, two observable CSTR screens exceeded the tolerance on new
data, and the certificate caught them. Certification is therefore not a formality but the step that detects
optimistic thresholds.
\item \emph{Retries are a different population.} A screen certified on first proposals must not accept retries
unchecked, because feedback changes the agent's internal state. The first-proposal-only rule turns this
observation into a one-line safeguard. It removed 95--99\% of the unsafe executions of the affected screens with
no measurable loss of recovery.
\item \emph{Savings depend on what replaces a skipped call.} An unchecked accept saves a full validation. An
unchecked rejection of a good proposal forces a regeneration. Screening therefore pays when a validation costs
more than the break-even $c^{\ast}$, which was 3--18\,s for four of five models. This is the regime that motivates
selective validation in the first place: high-fidelity or long-horizon simulators, scenario sets and multi-unit
plants, where the cost of validation grows while the cost of the agent stays fixed.
\end{enumerate}

\begin{figure}[!htbp]
\centering
% Deployment guide distilled from the results (each branch is backed by a result in Sections 6-7).
\resizebox{\textwidth}{!}{%
\begin{tikzpicture}[font=\footnotesize, >={Latex[length=2mm]},
  q/.style={draw=black!70, fill=black!4, rounded corners=2pt, text width=46mm, align=center, inner sep=4pt},
  a/.style={draw=black!70, rounded corners=2pt, text width=46mm, align=center, inner sep=4pt},
  lab/.style={font=\scriptsize, fill=white, inner sep=1pt},
]
\node[q] (q1) at (0,0) {Does a validator call cost more than the break-even $c^{\ast}$ (Eq.~\ref{eq:breakeven})?\\[-1pt]{\scriptsize estimate $c^{\ast}$ from a short logged pilot}};
\node[a] (n1) at (0,-2.7) {Validate every proposal.\\[-1pt]{\scriptsize screening cannot pay off; rejections cost regenerations}};
\node[q] (q2) at (6.2,0) {Are failures retrieval failures?\\[-1pt]{\scriptsize grounding AUROC on development data clearly above 0.5}};
\node[a, draw=cComb, fill=cComb!10] (a2) at (6.2,-2.7) {Combined screen:\\observables + internals + grounding\\[-1pt]{\scriptsize (path planning: certified gains up to +15 points)}};
\node[a, draw=cObs, fill=cObs!8] (a3) at (12.4,0) {Observable screen; internals add little.\\[-1pt]{\scriptsize where no plant model or plant data exist, an internals-only screen still rejects with 95--98\% precision}};
\draw[->] (q1) -- node[lab, left] {no} (n1);
\draw[->] (q1) -- node[lab, above] {yes} (q2);
\draw[->] (q2) -- node[lab, left] {yes} (a2);
\draw[->] (q2) -- node[lab, above] {no} (a3);
\node[draw=cRzero, fill=cRzero!12, rounded corners=2pt, text width=105mm, align=left, inner sep=5pt] (rules) at (9.3,-6.0)
 {\textbf{Always:} (i) set $t_{\mathrm{acc}}$ with a confidence bound, not a point estimate;
  (ii) certify on a sealed set large enough for Eq.~\eqref{eq:mmin};
  (iii) allow unchecked execution only for first proposals (R0);
  (iv) if the agent repairs proposals from the validator's feedback, an unchecked rejection trades recovery for a
  saved call: validate rejections when a recovered episode is worth more than the calls they save;
  (v) recalibrate and re-certify after a change in plant or fault population.};
\draw[->, black!60] (a2.south) -- (a2.south |- rules.north);
\draw[->, black!60] (a3.south) -- (a3.south |- rules.north);
\end{tikzpicture}}

\caption{Deployment guide derived from the results. The first question decides whether screening can pay off
(Section~\ref{sec:closedloop}). The second decides whether internal signals are worth adding to an observable
screen (Section~\ref{sec:res-rq3}). The rules at the bottom apply to every screen.}
\label{fig:guide}
\end{figure}

\subsection{Threats to validity}
\paragraph{Internal validity} All hypotheses, rules and amendments were pre-registered, and sealed partitions
were evaluated once. Two analyses are exploratory and labelled as such: the offline CSTR result of Qwen2.5-3B
and the closed-loop probes of that model. The first-proposal-only rule was introduced after observing the retry
failure. We therefore re-tested it on 400 new episodes (Section~\ref{sec:closedloop}). Teacher-forced re-runs
agreed with greedy generation for 92.8--99.9\% of answer tokens for the bfloat16 models but only 84--91\% for the
4-bit Qwen2.5-7B. For that model the internal features describe a closely related, not identical, decoding path.
This may contribute to its weaker gains from internals.

\paragraph{Construct validity} The CSTR deciding regions are domain-informed choices taken from the process
structure stated in the prompt. In a pre-registered sensitivity check we recomputed the CSTR grounding AUROC under
three alternative region definitions: the physics section, the section describing how each setpoint acts, and the
whole snapshot. For Qwen2.5-3B and Llama-3.2-3B no definition reached 0.55. For Llama, \emph{more} attention on
the physics section went with failure (0.35), as in the pre-registered H3 test. For Qwen2.5-7B, however, low
attention on the physics and action sections did predict failure (0.70 and 0.71), which our pre-registered
criterion had not anticipated. For that model the CSTR failures are partly failures to consult the physics, and
the retrieval--reasoning contrast is one of degree rather than kind. Attention is used as a predictive signal, not as an explanation. The observable
baseline in path planning includes a shortest-path length and an endpoint check, which are partial validator
computations. This makes the baseline strong and the gains from internals conservative.

\paragraph{External validity} The study uses open-weight models of 1.5--7 billion parameters, one of them
quantised, on a single laptop GPU. Larger models make fewer retrieval errors and may therefore benefit less from
grounding, and API-served models do not expose internal signals at all. The CSTR is one simulated plant with four
fault families and one prompt, and the plant is paused while the agent and the validator run. We make no claim
about decision latency under real-time constraints. The validator doubles as the plant in closed loop, so model
mismatch between the digital twin and a real plant is outside the scope of this study.

\paragraph{Statistical conclusion validity} Certification in the CSTR was limited by the size of the
certification partitions (Eq.~\ref{eq:mmin}), so the absence of certified CSTR screens reflects power as well as
risk. Confidence intervals are per comparison. The certification bounds are Bonferroni-corrected over signal sets
but not over models or tolerances. Compute figures are specific to our hardware, but the linear re-costing makes
their dependence on the validator explicit and transferable.
